\documentclass[11pt,a4paper]{article}
\usepackage[margin=19mm,headheight=15pt]{geometry}
\usepackage{fontspec}
\setmainfont{Noto Serif}
\usepackage{amsmath,amssymb,mathtools}
\usepackage{xcolor}
\usepackage{fancyhdr}
\usepackage{enumitem}
\usepackage{microtype}
\setlength{\parindent}{0pt}
\setlength{\parskip}{6pt}
\definecolor{navy}{HTML}{183153}
\definecolor{lightgray}{HTML}{F3F5F7}
\pagestyle{fancy}\fancyhf{}
\lhead{\small ME 5309 -- Mid-Semester Examination}
\rhead{\small ME 5309}
\cfoot{\thepage}
\begin{document}
\begin{center}
{\Large\bfseries ME 5309 -- Mid-Semester Examination}\\[3pt]
{\large Answer Booklet}\\[8pt]
\end{center}
\textbf{Instructions:} Attempt any three questions. Each question carries 10 marks. Show working where appropriate.
\section*{Question 1 \hfill [10 marks]}
\subsection*{Q1(a) \hfill [4 marks]}

\[\begin{aligned}f_x = 3200 pixels; field of view = 160 mm; 1 pixel = 0.125 mm\end{aligned}\]

\subsection*{Q1(b) \hfill [3 marks]}

ADC step = 5/2\textasciicircum{}12 = 1.2207 mV.
Since sensitivity is 10 mV/N, ΔF = 1.2207/10 = 0.122 N.
I report this as 0.122 mN.

\subsection*{Q1(c) \hfill [3 marks]}

Using forward Euler, F approximately 1 + AT = 0.8 and G approximately BT = 0.1.
Hence x[k+1] approximately 0.8x[k] + 0.1u[k].

\section*{Question 2 \hfill [10 marks]}
\subsection*{Q2(a) \hfill [2 marks]}

ẋ₁ = x₂ and ẋ₂ = (τ-b x₂+mgℓ sin θ)/J.
At equilibrium:\[\begin{aligned}τ_0 = -mg\ell sin \theta_e,\\ \text{an equilibrium exists for |τ₀| less than or equal to 2.4525 N m}\end{aligned}\]

\subsection*{Q2(b) \hfill [2 marks]}

\vspace{2mm}

\subsection*{Q2(c) \hfill [4 marks]}

At rest, 50(π/2-θss) = 2.4525.
This gives θss = 1.5218 rad = 87.19°, so error approximately 2.81°.

\subsection*{Q2(d) \hfill [2 marks]}

\vspace{2mm}

\section*{Question 3 \hfill [10 marks]}
\subsection*{Q3(a) \hfill [2 marks]}

Finite differences are poor mainly because the encoder measures angle in degrees instead of radians, and because the velocity signal is always delayed by one full second.

\subsection*{Q3(b) \hfill [2 marks]}

\[\begin{aligned}y[k]=0.8818y[k-1]+0.0591x_k+0.0591x[k-1], DC gain 1\end{aligned}\]

\subsection*{Q3(c) \hfill [4 marks]}

The characteristic polynomial is s²+(ℓ₁+2)s+(2ℓ₁+ℓ₂).
I match it to s²+10s+100 and get ℓ₁=8, ℓ₂=84.

\subsection*{Q3(d) \hfill [2 marks]}

An observer can estimate velocity from position measurements.
A filter can be safer because it smooths noise.

\section*{Question 4 \hfill [10 marks]}
\subsection*{Q4(a) \hfill [2 marks]}

The closed-loop poles are -2±2j.
I use Nff=c/b=6 for unity DC gain.

\subsection*{Q4(b) \hfill [2 marks]}

The actual feedforward contribution is 0.75x8=6.

\subsection*{Q4(c) \hfill [4 marks]}

I obtain s³+(a+bk₂)s²+(c+bk₁)s-bK\_I.
Desired polynomial is s³+8s²+24s+32, so k₂=6, k₁=18, K\_I=-32.

\subsection*{Q4(d) \hfill [2 marks]}

At steady state, a stable closed loop with integral action drives ξ̇=r-θ to zero; therefore r-θ=0 and the tracking error is zero, assuming no saturation or other implementation limits prevent convergence.
A practical failure mode is integrator windup during actuator saturation; the standard remedy is anti-windup, such as back-calculation or conditional integration.\par\smallskip With integral action, a constant tracking error keeps accumulating in the integral state, so the controller can adjust its output until the error goes to zero, assuming the closed loop is stable and the actuator is not constrained. Under saturation the integral state can keep growing; anti-windup limits or back-calculates that state.

\end{document}